We propose intent-aware retrieval and semantics-prethat preserves the semantic integrity of the retrieved evidence to alleviate the problem of incomplete information retrieval in conventional RAG systems.
We utilize SLM as our core retriever for fast and proper instruction following capability.
Our work yields competitive results against state-of-the-art RAG systems at a much lower latency.
These results broaden the exploration of efficient semantic-aware RAG for LLM.
%Our work expands the exploration of an efficient and semantically complete RAG system with SLM.% at its heart.

%This paper proposes a retrieval-enhanced generation framework based on small language model, SLM-RAG, which aims to achieve efficient, credible and interpretable semantic retrieval and knowledge integration under resource-constrained conditions. The framework significantly improves retrieval recall, evidence integrity, and contextual consistency while maintaining lightweight by introducing mechanisms such as hypothetical query reasoning (HQ Reasoning), hybrid search fusion (BM25 combined with Dense), forensic completion, coverage physical examination, and structured knowledge integration. Compared with traditional RAG systems that rely on large-scale language models, SLM-RAG shows obvious advantages in computational efficiency and response speed, while having stronger controllability and interpretability in terms of multi-hop correlation and evidence traceability.
%Experimental results show that SLM-RAG achieves stable performance improvement in multi-domain Q&A and enterprise knowledge retrieval tasks, which verifies the potential value of small language models in complex retrieval and inference tasks. The SLM-3-driven knowledge integration module in the framework transforms search results into structured, contextually consistent knowledge outputs, providing reliable semantic support for upper-level Q&A, report generation, or agent decision-making. This design idea of "lightweight autonomy + hierarchical collaboration" provides new possibilities for building efficient and low-cost cognitive information systems in the future.
%However, current research still has certain limitations. First, the collaboration between SLM modules still relies on artificially set stage boundaries, and end-to-end adaptive optimization has not yet been realized. Secondly, the strategy of hybrid retrieval still needs to be further optimized under large-scale heterogeneous corpus to avoid redundant overlap between sparse and dense retrieval. In addition, although the knowledge integration process of SLM-3 improves semantic consistency, it may still be limited by the memory and context window length of the small model in long-chain multi-hop inference tasks. Future work will focus on introducing reinforcement learning and dynamic programming mechanisms to achieve cross-module self-scheduling and continuous optimization, while exploring the expansion potential of SLM-RAG in multimodal retrieval and cross-language knowledge integration.
%Overall, SLM-RAG demonstrates a new paradigm of retrieval enhancement that combines lightweight, autonomy, and interpretability, laying an important foundation for promoting the application of small language models in the fields of complex information retrieval, knowledge question answering, and intelligent agents.